\documentclass[../main.tex]{subfiles}
\begin{document}

\textit{This chapter presents the implemented system through screenshots taken from the running
application with the interface in English, then the testing strategy and its results, and finally
the deployment.}

\section{Development environment}
\begin{table}[H]
\centering
\caption{Development environment and project size}\label{tab:env}
\renewcommand{\arraystretch}{1.2}
\begin{tabular}{|p{5cm}|p{8.7cm}|}
\hline
Operating system & Ubuntu Linux, Docker Engine with Compose v2 \\ \hline
Frontend & Node.js 22, pnpm 10, Vite, Vue 3, TypeScript \\ \hline
Backend & Python 3.12, FastAPI, SQLAlchemy 2, Alembic \\ \hline
Database & MySQL 8.4 (production), SQLite (development, tests) \\ \hline
Source code & about 13,300 lines frontend, 3,400 lines backend (non-blank) \\ \hline
Test code & about 5,200 lines frontend, 2,800 lines backend \\ \hline
Repository & 65 commits; 26 Vue components; 51 API endpoints; 10 tables; 9 migrations \\ \hline
\end{tabular}
\end{table}

\section{Implemented features}
\subsection{Landing page and authentication}
A visitor first sees a marketing page with the features, a demo video and pricing; the sign-in
button is at the top right and the language can be switched between English and Vietnamese
(Figure~\ref{fig:landing.png}). The demo video is a genuine screen recording: a script drives the
real application in a browser and captures it, so it can only show features that work.
Registration asks for a name, an email and a password---never a role (Figure~\ref{fig:login.png}).

\shot{landing.png}{0.95}{Landing page}
\shot{login.png}{0.95}{Sign-in dialog}

\subsection{Projects, members and roles}
The dashboard lists the user's projects with the user's role in each, creates projects and manages
the team (Figure~\ref{fig:dashboard.png}). Adding a member by email sends an invitation; the person
joins only after accepting it.

\shot{dashboard.png}{0.95}{Project dashboard with per-project roles}

\subsection{Kanban board and table view}
The board has four columns. Cards show the key, title, priority, assignees' avatars, the deadline
(``in 9d'', ``3d overdue''), the blocking reason and a flame icon for tasks on the critical path.
Ordering is deadline-first: a low-priority task due yesterday ranks above a critical one due next
month. The same data is available as a table (Figures~\ref{fig:board-kanban.png} and
\ref{fig:board-table.png}); the key \code{b} switches between them.

\shot{board-kanban.png}{0.95}{Kanban board}
\shot{board-table.png}{0.95}{Table view}

Every action has a key (Figure~\ref{fig:shortcuts.png}): $h/j/k/l$ move the selection, Space cycles
the status, $H/L$ move a task between columns, 1--4 set the priority, $n$ creates a task and $/$
searches.

\shot{shortcuts.png}{0.8}{Keyboard shortcuts}

\subsection{Task inspector: discussion, evidence and history}
Opening a task shows its fields, dependencies and acceptance criteria, a discussion thread with
@mentions and one-level replies, attachments, and the change history recorded by the server
(Figure~\ref{fig:task-detail.png}). Moving a task to DONE invites the user to attach evidence that
the work is finished.

\shot{task-detail.png}{0.95}{Task inspector with discussion and history}

\subsection{Scrum}
The Scrum panel has five tabs. \emph{Sprints} (Figure~\ref{fig:scrum-sprints.png}) lists planned,
active and completed sprints with their points and commitment; a PM creates, starts and completes
them, choosing where unfinished work goes. \emph{Planning} (Figure~\ref{fig:scrum-planning.png})
moves items between the product backlog and a sprint. \emph{Burndown}
(Figure~\ref{fig:scrum-burndown.png}) plots remaining points against the ideal line;
\emph{Velocity} (Figure~\ref{fig:scrum-velocity.png}) compares committed and completed points per
sprint with the average; \emph{Retrospective} (Figure~\ref{fig:scrum-retro.png}) asks the AI for
what went well, what to improve and concrete actions, grounded in the sprint's numbers.

\shot{scrum-sprints.png}{0.95}{Scrum panel: sprint lifecycle}
\shot{scrum-planning.png}{0.95}{Scrum panel: backlog planning}
\shot{scrum-burndown.png}{0.95}{Scrum panel: burndown chart}
\shot{scrum-velocity.png}{0.95}{Scrum panel: velocity}
\shot{scrum-retro.png}{0.95}{Scrum panel: AI retrospective}

\subsection{AI tools}
Four further AI tools are available from the \emph{AI Tools} menu. The weekly summary reports
progress, risks and priorities (Figure~\ref{fig:ai-weekly.png}); the meeting-minutes generator
turns raw notes into topics, action items with an owner and priority, and decisions
(Figure~\ref{fig:ai-minutes.png}); smart assignment shows each member's load and skills and
recommends an assignee with a rationale and a risk assessment (Figure~\ref{fig:ai-assign.png}); the
goal decomposer breaks a goal into dependent subtasks that can be inserted into the board in one
step (Figure~\ref{fig:ai-decomposer.png}).

\shot{ai-weekly.png}{0.95}{AI weekly progress summary}
\shot{ai-minutes.png}{0.95}{AI meeting minutes}
\shot{ai-assign.png}{0.95}{AI workload analysis and assignment}
\shot{ai-decomposer.png}{0.95}{AI goal decomposer}

\subsection{Dependency graph, notifications and profile}
The DAG view lists tasks in topological order and marks the critical path
(Figure~\ref{fig:dag.png}). The notification feed reports mentions, replies, assignments and
invitations---never the user's own actions (Figure~\ref{fig:notifications.png}). The profile page
edits the name, skills (used by the assignment recommender) and avatar (Figure~\ref{fig:profile.png}).

\shot{dag.png}{0.95}{DAG execution order and critical path}
\shot{notifications.png}{0.95}{Notification feed}
\shot{profile.png}{0.95}{Profile page}

\section{Testing}
\subsection{Strategy}
Testing had four layers:
\begin{enumerate}
  \item \textbf{Backend tests} (pytest, 13 files) exercise the API through a test client: every
        role and permission rule (including the negative cases---a member receiving 403, an outsider
        receiving 404), the status cycle, Scrum rules, history, migrations, startup safety checks,
        and the AI cascade with each tier simulated at the HTTP boundary.
  \item \textbf{Frontend tests} (Vitest, 19 files) cover the pure algorithms, both stores, the
        keyboard dispatcher and the components, mounted in a simulated DOM. One test reads every
        component template and fails if any literal text is not translated.
  \item \textbf{End-to-end smoke test}: a script seeds a running instance through the public API
        and performs 32 checks across all features, including the four AI endpoints against the
        real model.
  \item \textbf{Browser automation}: Playwright drives the real application to record the demo
        video and the screenshots in this chapter, and scans each screen for text in the wrong
        language.
\end{enumerate}
Tests were themselves tested by mutation: a guard was removed or a value changed on purpose, and the
change was accepted only if a test failed. For example, removing the check that a sprint belongs to
the task's project made two Scrum tests fail; removing the filter on accepted members made the
assignment test fail.

\subsection{Results}
\begin{table}[H]
\centering
\caption{Test results}\label{tab:tests}
\renewcommand{\arraystretch}{1.2}
\begin{tabular}{|p{6cm}|c|c|}
\hline
\textbf{Suite} & \textbf{Tests} & \textbf{Result} \\ \hline
Backend (pytest, SQLite) & 177 & all passed \\ \hline
Frontend (Vitest) & 442 & all passed \\ \hline
TypeScript type check and production build & -- & passed \\ \hline
End-to-end smoke test on the public instance & 32 checks & 0 failed \\ \hline
Mixed-language scan of 20 screens (English UI) & 20 & none found \\ \hline
\end{tabular}
\end{table}

\subsection{Defects found}
Testing found real defects, each recorded in the development book with its cause. A selection:
\begin{longtable}{|p{1.2cm}|p{9cm}|p{3.3cm}|}
\caption{Selected defects found and fixed}\label{tab:defects}\\
\hline
\textbf{ID} & \textbf{Defect} & \textbf{How found} \\ \hline
\endfirsthead
\hline
\textbf{ID} & \textbf{Defect} & \textbf{How found} \\ \hline
\endhead
F-32 & The published backend image contained the developer's \code{.env} with the signing secret and a database. & Inspecting the built image \\ \hline
F-37 & The documented deploy command copied development settings and databases over production. & Reviewing the command before running it \\ \hline
F-45 & All attachments and avatars rendered as broken images: an \code{<img>} tag cannot send an authorisation header. & Manual testing \\ \hline
F-47 & SQLite foreign keys were never enforced, so deleting a task left orphaned rows. & Migration to MySQL \\ \hline
F-48 & A task could be planned into another project's sprint. & Writing the Scrum tests \\ \hline
F-49 & A finished task was labelled ``10d overdue''. & Screenshot review \\ \hline
F-50 & The ``AI'' goal decomposer never called a model; it returned a fixed template. & Testing against the real model \\ \hline
F-51 & The assignment recommender could suggest someone who had not yet accepted the invitation. & Code review while testing AI \\ \hline
F-53 & A proposed change faked the PM role in the browser after the server refused it. & Review of incoming changes \\ \hline
-- & 21 of 26 components contained untranslated text, so the two languages were mixed. & Team member's report; fixed and guarded by a test \\ \hline
\end{longtable}

\section{Deployment}
The production configuration runs four containers: MySQL, the FastAPI backend, nginx with the
built frontend, and \code{cloudflared}. A deployment script builds the images, verifies that no
secret or database is inside them, backs up the database, applies migrations in a separate step,
starts the services and checks that the instance is healthy, at the latest schema revision and in
a production posture (no demo accounts, no wildcard CORS, a non-default secret). The tunnel
publishes the instance at \url{https://koshi.huynhca.com}, where the team tested it with three
seeded accounts. Continuous integration on GitHub Actions runs both test suites on every push to the main and development branches.

\end{document}
